\id{IRSTI 61.67.81}{}

\begin{header}
\swa{N.V.~Seraya, V.V.~Litvinov, G.K.~Daumova, D.E.~Sovetkhanov, M.~Rutkowska-Gorczyca, M.M.~Baimurzina}{PHYSICO-MECHANICAL PROPERTIES AND DURABILITY OF GEOCOMPOSITE BASED ON LOW-PRESSURE POLYETHYLENE AND GEOTEXTILE DURING LONG-TERM OPERATION}

\tsp{1}N.V.~Seraya,
\tsp{2}V.V.~Litvinov,
\tsp{1}G.K.~Daumova\envelope,
\tsp{1}D.E.~Sovetkhanov,
\tsp{3}M.~Rutkowska-Gorczyca,
\tsp{1}M.M.~Baimurzina
\end{header}

\begin{affil}
\tsp{1}D.Serikbayev East Kazakhstan Technical University, Ust-Kamenogorsk, Kazakhstan,

\tsp{2}Proektno-ekologicheskoe bjuro LLP, Ust-Kamenogorsk, Kazakhstan,

\tsp{3}Wrocław University of Science and Technology, Wrocław, Poland

\corrauthor{Correspondent-author: gulzhan.daumova@mail.ru}
\end{affil}

The paper presents the results of experimental studies of the
physico-mechanical properties of geocom\-posite used as a covering
material for winter blocks of TMF dumps of the Kounrad mine. The changes
in the basic strength and deformation characteristics of the material
depending on the service life (1 month, 2 years, 3 years) have been
studied. The tests were carried out according to GOST standards with the
determination of surface density, tensile and tear strength, elongation,
yield strength, resistance to aggressive media and repeated freezing and
thawing. The results showed that with increasing service life, there is
an increase in surface density (from 445 to 604 g/m\tsp{2})
and tensile strength (up to 9.5 kN/m), while reducing elongation.
Geocomposite demonstrates high resistance to aggressive environments (up
to 96\%) and frost resistance, which confirms its durability and
suitability for long-term use under conditions of variable climatic
loads.

{\bfseries Keywords:} geocomposite, low-pressure polyethylene, geotextile,
physical and mechanical properties, strength, deformability, frost
resistance, resistance to aggressive environments, operation,
durability.

\begin{header}
ФИЗИКО-МЕХАНИЧЕСКИЕ СВОЙСТВА И ДОЛГОВЕЧНОСТЬ ГЕОКОМПОЗИТА НА ОСНОВЕ ПОЛИЭТИЛЕНА НИЗКОГО ДАВЛЕНИЯ И ГЕОТЕКСТИЛЯ ПРИ ДЛИТЕЛЬНОЙ ЭКСПЛУАТАЦИИ

\tsp{1}Н.В.~Серая,
\tsp{2}В.В.~Литвинов,
\tsp{1}Г.К.~Даумова\envelope,
\tsp{1}Д.Е.~Советханов,
\tsp{3}M.~Rutkowska-Gorczyca,
\tsp{1}М.М.~Баймурзина
\end{header}

\begin{affil}
\tsp{1}Восточно-Казахстанский технический университет им. Д.Серикбаева, Усть-Каменогорск, Казахстан,

\tsp{2}ТОО «Проектно-экологическое бюро», Усть-Каменогорск, Казахстан,

\tsp{3}Wrocław University of Science and Technology, Вроцлав, Польша,

e-mail: gulzhan.daumova@mail.ru
\end{affil}

В работе представлены результаты экспериментальных исследований
физико-механических сво\-йств геокомпозита, применяемого в качестве
укрывного материала зимних блоков отвалов ТМО Коунрадского рудника.
Изучены изменения основных прочностных и деформационных характеристик
материала в зависимости от срока эксплуатации (1 месяц, 2 года, 3 года).
Испытания проводились по стандартам ГОСТ с определением показателей
поверхностной плотности, прочности при растяжении и разрыве,
относительного удлинения, предела текучести, устойчивости к агрессивным
средам и многократному замораживанию-оттаиванию. Результаты показали,
что с увеличением срока эксплуатации наблюдается рост поверхностной
плотности (с 445 до 604 г/м²) и прочности при растяжении (до 9,5 кН/м),
при одновременном снижении относительного удлинения. Геокомпозит
демонстрирует высокую устойчивость к агрессивным средам (до 96 \%) и
морозостойкость, что подтверждает его долговечность и пригодность для
длительного использования в условиях переменных климатических нагрузок.

{\bfseries Ключевые слова:} геокомпозит, полиэтилен низкого давления,
геотекстиль, физико-механичес\-кие свойства, прочность, деформативность,
морозостойкость, устойчивость к агрессивным средам, эксплуатация,
долговечность.

\begin{header}
ҰЗАҚ МЕРЗІМДІ ПАЙДАЛАНУ КЕЗІНДЕ ТӨМЕН ҚЫСЫМДЫ ПОЛИЭТИЛЕН ЖӘНЕ ГЕОТЕКСТИЛЬ НЕГІЗІНДЕГІ ГЕОКОМПОЗИТТІҢ ФИЗИКАЛЫҚ-МЕХАНИКАЛЫҚ ҚАСИЕТТЕРІ МЕН БЕРІКТІГІ

\tsp{1}Н.В.~Серая,
\tsp{2}В.В.~Литвинов,
\tsp{1}Г.К.~Даумова\envelope,
\tsp{1}Д.Е.~Советханов,
\tsp{3}M.~Rutkowska-Gorczyca,
\tsp{1}М.М.~Баймурзина
\end{header}

\begin{affil}
\tsp{1}Д. Серікбаев атындағы Шығыс Қазақстан техникалық университеті, Өскемен, Қазақстан,

\tsp{2}"Проектно-экологическое бюро" ЖШС, Өскемен, Қазақстан,

\tsp{3}Wrocław University of Science and Technology, Вроцлав, Польша,

e-mail: gulzhan.daumova@mail.ru
\end{affil}

Жұмыста Қоңырат кенішінің ТМО үйінділерінің қысқы блоктарының жабын
материалы ретінде қолданылатын геокомпозиттің физика-механикалық
қасиеттерін эксперименттік зерттеу нәтижелері келтірілген. Пайдалану
мерзіміне (1 ай, 2 жыл, 3 жыл) байланысты материалдың негізгі беріктігі
мен деформациялық сипаттамаларының өзгеруі зерттелді. Сынақтар беттің
тығыздығы, созылу және үзілу беріктігі, салыстырмалы ұзарту, аққыштық
шегі, агрессивті ортаға төзімділік және бірнеше рет мұздату-еріту
көрсеткіштерін анықтай отырып, МЕСТ стандарттары бойынша жүргізілді.
Нәтижелер пайдалану мерзімінің ұлғаюымен салыстырмалы ұзарудың бір
мезгілде төмендеуімен беттің тығыздығының (445-тен 604
г/м\tsp{2}-ге дейін) және созылу кезіндегі беріктіктің (9,5
кН/м дейін) өсуі байқалатынын көрсетті. Геокомпозит агрессивті ортаға
жоғары төзімділікті (96\% дейін) және аязға төзімділікті көрсетеді, бұл
оның беріктігі мен өзгермелі климаттық жүктемелер жағдайында ұзақ
мерзімді пайдалануға жарамдылығын растайды.

{\bfseries Түйін сөздер:} геокомпозит, төмен қысымды полиэтилен,
геотекстиль, физикалық-механикалық қасиеттер, беріктік, деформация,
аязға төзімділік, агрессивті ортаға төзімділік, эксплуатация, ұзақ
мерзімді пайдалану.

\begin{multicols}{2}
{\bfseries Introduction.} Modern mining technologies are increasingly
introducing geosynthetic materials to solve the problems of engineering
protection and stabilization of landfills, tailings dumps and other
man-made structures. Of particular importance in this context are
geocomposites - multilayer materials that combine the functions of
sealing, reinforcement, thermal insulation and protection against
erosion processes {[}1-3{]}. Due to the combination of polymer and
textile components, geocomposites provide high performance properties
and durability with a relatively small weight.

One of the most common types of such materials is geocomposite based on
low-pressure polyethylene (LPP) and non-woven geotextiles. This type of
geomaterial is actively used as a covering and heat-protective coating
in the conservation of winter blocks of landfills at mining enterprises.
However, during long-term operation, under the influence of temperature
fluctuations, freeze-thaw cycles, solar radiation and mechanical loads,
physico-chemical and morphological changes occur in the polymer
structure, affecting its strength and deformability {[}4, 5{]}.

Authors {[}6, 7{]} showed that over time, the main physical and
mechanical parameters of geosynthetics, such as tensile strength, yield
strength and elongation, change depending on the degree of aging of the
material. Polyethylene and polypropylene components are characterized by
a gradual increase in stiffness and density while decreasing plasticity.
This is due to the processes of relaxation of internal stresses,
accumulation of microdefects and compaction of the material structure.

Special attention is paid to the behavior of geosynthetics when exposed
to aggressive chemical environments and freeze-thaw cycles, since these
factors determine their service life in real-world operating conditions
{[}8,9{]}. It is known that for geocomposites used in mining
enterprises, high frost resistance and chemical inertia are key
parameters of durability.

Despite a significant number of studies devoted to microstructural
changes in geosynthetics, the relationship between density, strength
characteristics and deformation properties of materials during long-term
operations remains insufficiently studied. This is especially true for
geocomposites used as thermal protective coatings for landfills in a
sharply continental climate.

The purpose of this work is an experimental study of the
physico-mechanical properties of a LPP-based geocomposite and geotextile
used to shelter winter blocks of TMF dumps at the Kounrad mine, and an
assessment of changes in its strength and deformation characteristics
depending on the service life.

{\bfseries Materials and methods.} \emph{The object of research.} The
object of the study was a geocomposite based on low-pressure
polyethylene (LPP) and non-woven geotextile, used as a heat-protective
coating for winter blocks of landfills of technogenic mineral formations
(TMF) of the Kounrad mine. The material is a two-layer structure that
combines the moisture resistance of the polymer layer and the mechanical
strength of the textile fabric.

Geocomposite samples that had been in use for various periods of time
were selected for testing:

- 1 month - shelter 2022 (block 26, dump 16);

- 2 years - shelter in 2021 (block 20, dump 16);

- 3 years - shelter in 2020 (block 3, dump 16).

This approach made it possible to evaluate the effect of the duration of
operation on the physical and mechanical properties of the material.

\emph{Test procedure}. The tests of the physical and mechanical
properties of the geocomposite were carried out in accordance with the
requirements of the current national standards - GOST R 50277, 50300-92,
56586, 11262-2017, 55035-2012, 55032-2012 {[}10-15{]}.

The following parameters were determined for each sample:

- surface density, g/m2 (GOST R 50277) {[}10{]};

- tensile strength in the longitudinal and transverse directions, kN/m
(GOST R 50300-92) {[}11{]};

- elongation at maximum load, \% (GOST R 50300-92) {[}11{]};

- tensile strength and elongation at break, kN/m and \% (GOST R 56586
{[}12{]}, GOST 11262-2017) {[}13{]};

- yield strength and elongation at yield strength (GOST R 56586
{[}12{]}, GOST 11262-2017) {[}13{]};

- resistance to aggressive environments (GOST R 55035-2012) {[}14{]};

- resistance to repeated freezing and thawing --- frost resistance (GOST
R 55032-2012) {[}15{]}.

Tensile tests were carried out on a universal testing machine at a
constant rate of deformation, with maximum load and elongation fixed
until the sample collapses. Each test was performed in three parallel
repeats, and the results were averaged with a standard deviation.

To determine chemical resistance, the samples were exposed to acidic (pH
≈ 2) and alkaline (pH ≈ 12) media at a temperature of (60 ± 1) °C for 72
hours. Upon completion of the tests, the materials were rinsed with
running water, rinsed with distilled water and dried for at least 24
hours at room temperature. Next, the residual tensile strength was
determined. The degree of resistance to aggressive media was assessed by
the ratio of the strength of the treated samples to the strength of
control samples that were not exposed to reagents.

The assessment of resistance to multiple cycles of freezing and thawing
was carried out by comparing the mechanical characteristics of control
samples and samples that had undergone 10-20 cycles of freezing (at -20
°C) and thawing (at +20 °C) in an aqueous environment. At the end of the
cycle, the samples were subjected to tensile tests. Frost resistance was
expressed as a percentage of strength retention after cyclic exposure.

The measurement results were processed statistically, with the
calculation of average values, standard deviations and the construction
of comparative dependencies between the service life and changes in
physical and mechanical characteristics. Grapho-analytical methods were
used to interpret relationships.

A correlation analysis was performed to identify the relationships
between the physical and mechanical characteristics of the geocomposite.
The method is based on the calculation of Pearson paired linear
correlation coefficients, which allow quantifying the degree and
direction of the relationship between two random variables. The
correlation coefficient r\tsb{xy} was calculated using the
formula (1):

\begin{equation}
r_{xy} = \frac{\sum_{i = 1}^{n}{\left( x_{i} - \overline{x} \right) \cdot \left( y_{i} - \overline{y} \right)}}{\sqrt{\sum_{i = 1}^{n}{\left( x_{i} - \overline{x} \right)^{2} \cdot}\sum_{i = 1}^{n}\left( y_{i} - \overline{y} \right)^{2}}}
\end{equation}

where \(x_i\), \(y_i\) are the values of the
parameters being compared;

\(\overline{x},\overline{y}\) - their average values;

\(n\) is the number of observations.

The values of the r\tsb{xy} coefficient range from -1 to +1.
Positive values correspond to a direct correlation (an increase in one
parameter is accompanied by an increase in the other), negative values
correspond to an inverse relationship, and values close to zero indicate
the absence of a relationship. The average values of the parameters
obtained from the results of three series of tests of geocomposite
samples with different service life (1 month, 2 years, 3 years) were
used for the analysis.

{\bfseries Results and discussion.} The results of experimental studies of
the physico-mechanical properties of the geocomposite are presented in
Table 1. Analysis of the data obtained made it possible to establish
patterns of changes in the strength and deformation characteristics of
the material depending on its service life as a heat-protective coating
for winter blocks of landfills.

One of the key indicators of the state of a geosynthetic material is its
surface density, which determines the specific gravity of the web and
reflects the degree of its uniformity. According to the experimental
data obtained (Table 1), the surface density of the geocomposite
increases from 445 g/m\tsp{2} (sample 2022, 1 month of
operation) to 604 g/m\tsp{2} (sample 2020, 3 years of
operation). The increase in this indicator is associated with the
penetration of mineral and organic pollutants into the structure of the
material both from the outside (surface infiltration) and from the
inside --- from the side of contact with the soil. An increase in
density leads to an increase in the strength, rigidity and wear
resistance of the material, but it is accompanied by a decrease in
breathability and stretchability. This indicates a gradual consolidation
of the geocomposite structure during operation, which was also noted by
other researchers {[}2, 16{]}.

Tensile strength is the main characteristic determining the bearing
capacity of geotextile and polymer components. According to the test
results, the strength of the geocomposite increases longitudinally from
8.6 to 9.5 kN/m, and in the transverse direction from 6.3 to 7.7 kN/m.
The increase in strength indicators confirms the effect of "compaction"
of the material during long-term operation, which is associated with a
decrease in fiber mobility and increased adhesion between the layers of
LPP and geotextile. At the same time, an increase in strength is
accompanied by an increase in stiffness, which is consistent with the
literature data on the behavior of thermally bonded geosynthetics after
cycles of thermo-mechanical aging {[}17{]}.

Elongation characterizes the ductility and deformability of a material.
As the service life of the geocomposite increases, the elongation index
decreases from 66 to 44\% in the longitudinal direction and from 67 to
43\% in the transverse direction (Table 1). The decrease in plasticity
is explained by an increase in the density and rigidity of the material
due to the compaction of the structure and partial aging of the polymer.
This is a typical pattern for polyethylene and polypropylene
geosynthetics exposed to atmospheric aging {[}18{]}.

The tensile strength increases from 1.74 to 2.49 kN/m in the
longitudinal direction and from 1.85 to 2.95 kN/m in the transverse
direction (Table 1). This indicates that geocomposite retains high
strength properties even after prolonged exposure to natural factors. It
is likely that partial hardening of the material occurs due to
structural "self-sealing" and stress redistribution in the fibers during
operation. A similar effect is noted in studies on the durability of
geosynthetic materials in the climatic conditions of Central Asia
{[}19{]}.
\end{multicols}

\tcap{Table 1 - Basic physical and mechanical characteristics of geocomposite samples}
\begin{longtblr}[
  label = none,
  entry = none,
]{
  width = \linewidth,
  colspec = {Q[20]Q[600]Q[140]Q[140]Q[140]},
  cells = {font = \footnotesize},
  row{1} = {c,font=\bfseries\footnotesize},
  row{2} = {c,font=\bfseries\footnotesize},
  cell{1}{1} = {r=2}{},
  cell{1}{2} = {r=2}{},
  cell{1}{3} = {c=3}{},
  cell{3}{1} = {c},
  cell{3}{3} = {c},
  cell{3}{4} = {c},
  cell{3}{5} = {c},
  cell{4}{1} = {c},
  cell{4}{3} = {c},
  cell{4}{4} = {c},
  cell{4}{5} = {c},
  cell{5}{1} = {c},
  cell{5}{3} = {c},
  cell{5}{4} = {c},
  cell{5}{5} = {c},
  cell{6}{1} = {c},
  cell{6}{3} = {c},
  cell{6}{4} = {c},
  cell{6}{5} = {c},
  cell{7}{1} = {c},
  cell{7}{3} = {c},
  cell{7}{4} = {c},
  cell{7}{5} = {c},
  cell{8}{1} = {c},
  cell{8}{3} = {c},
  cell{8}{4} = {c},
  cell{8}{5} = {c},
  cell{9}{1} = {c},
  cell{9}{3} = {c},
  cell{9}{4} = {c},
  cell{9}{5} = {c},
  cell{10}{1} = {c},
  cell{10}{3} = {c},
  cell{10}{4} = {c},
  cell{10}{5} = {c},
  cell{11}{1} = {c},
  cell{11}{3} = {c},
  cell{11}{4} = {c},
  cell{11}{5} = {c},
  cell{12}{1} = {c},
  cell{12}{3} = {c},
  cell{12}{4} = {c},
  cell{12}{5} = {c},
  cell{13}{1} = {c},
  cell{13}{3} = {c},
  cell{13}{4} = {c},
  cell{13}{5} = {c},
  cell{14}{1} = {c},
  cell{14}{3} = {c},
  cell{14}{4} = {c},
  cell{14}{5} = {c},
  cell{15}{1} = {c},
  cell{15}{3} = {c},
  cell{15}{4} = {c},
  cell{15}{5} = {c},
  cell{16}{1} = {c},
  cell{16}{3} = {c},
  cell{16}{4} = {c},
  cell{16}{5} = {c},
  cell{17}{1} = {c},
  cell{17}{3} = {c},
  cell{17}{4} = {c},
  cell{17}{5} = {c},
  hlines,
  vlines,
}
№ & Parameter & Geocomposite sample & & \\
 & & 26th block, 16th dump, 2022 y (1 month) & 20th block, 16th dump, 2021 y, (2 years) & 3rd block, 16th dump, 2020 y (3 years old) \\
1 & Surface density, g/m\textsuperscript{2} (according to GOST R 50277) & 445±16 & 523±41 & 604±77 \\
2 & Tensile strength (longitudinal direction), kN/m (according to GOST R 50300-92) [11] & 8,6±0,5 & 8,9±0,3 & 9,5±0,8 \\
3 & Tensile strength (transverse direction), kN/m (according to GOST R 50300-92) [11] & 6,3±0,4 & 7,6±0,4 & 7,7±0,9 \\
4 & Elongation at maximum load (in the longitudinal direction), \% (according to GOST R 50300-92) [11] & 66±4 & 36±2 & 44±2 \\
5 & Elongation at maximum load (in the transverse direction), \% (according to GOST R 50300-92) [11] & 67±5 & 53±2 & 43±3 \\
6 & Tensile strength (in the longitudinal direction), kN/m (according to GOST R 56586 [12], GOST 11262-2017 [13]) & 1,74±0,13 & 2,46±0,59 & 2,49±0,29 \\
7 & Tensile strength (in the transverse direction), kN/m (according to GOST R 56586 [12], GOST 11262-2017 [13]) & 1,85±0,17 & 2,87±0,59 & 2,95±0,25 \\
8 & Elongation at break (in the longitudinal direction), \% (according to GOST R 56586 [12], GOST 11262-2017 [13]) & 234±34 & 156±30 & 93±6 \\
9 & Elongation at break (in the transverse direction), \% (according to GOST R 56586 [12], GOST 11262-2017 [13]) & 94±23 & 177±25 & 145±20 \\
10 & Yield strength (in the longitudinal direction), kN/m (according to GOST R 56586 [12], GOST 11262-2017 [13]) & 2,13±0,21 & 3,62±0,30 & 2,69±0,29 \\
11 & Yield strength (in the transverse direction), kN/m (according to GOST R 56586 [12], GOST 11262-2017 [13]) & 2,23±0,15 & 3,64±0,31 & 3,14±0,21 \\
12 & Elongation at yield strength (in the longitudinal direction), \% (according to GOST R 56586 [12], GOST 11262-2017[13]) & 11±1 & 15±1 & 16±2 \\
13 & Elongation at yield strength (in the transverse direction), \% (according to GOST R 56586 [12], GOST 11262-2017 [13]) & 12±1 & 13±1 & 12±1 \\
14 & Resistance to aggressive environments, \% (according to GOST R 55035 [14]) & 92,1 & 93,4 & 94,9 \\
15 & Resistance to repeated freezing and thawing (frost resistance), \% (according to GOST R 55032-2012 [15]) & 92,2 & 95,5 & 93,1 
\end{longtblr}

\begin{multicols}{2}
The change in elongation at break is multidirectional. In the
longitudinal direction, the elongation decreases from 234 to 93\%,
reflecting a decrease in the ductility and elastic properties of the
material; in the transverse direction, on the contrary, it increases
from 94 to 145\%. This anisotropy is due to the design features of the
geotextile web: the longitudinal direction is more durable but less
elastic, while the transverse direction retains the ability to deform
under load, compensating for tensile forces.

Yield strength indicators reflect the ability of a material to maintain
its shape under prolonged loads. In the longitudinal direction, the
yield strength increases from 2.13 to 2.69 kN/m, in the transverse
direction - from 2.23 to 3.14 kN/m. An increase in values indicates the
formation of a more stable polymer structure showing signs of
self-hardening. This is consistent with the known relaxation and creep
effects for polyethylene described in the papers {[}20, 21{]}.

The elongation at yield strength increases longitudinally from 11 to
16\% and remains almost unchanged in the transverse direction (12-13\%).
This data confirm that the material is able to partially retain its
deformability during the transition to the stage of plastic flow, which
is important for assessing its operational reliability.

According to the test results, the geocomposite demonstrates high
chemical resistance: resistance to acidic and alkaline environments is
92-96\%. After three years of operation, the material even increases its
stability, which is explained by the compaction of the structure and a
decrease in the permeability of fibers to reagents. These results
confirm the chemical inertia of LPP and polypropylene fibers {[}22{]}.

All the samples studied showed high resistance to repeated cycles of
freezing and thawing - the strength retention was 92-95\% (Table 1). In
some cases, even a slight increase in strength was observed, which is
associated with the effect of "shrinkage" and compaction of the
structure during freezing. These results confirm the possibility of
long-term use of geocomposite in a sharply continental climate without
loss of operational properties.

A comparison of the data obtained shows that with increasing service
life, there is a tendency to increase the strength characteristics and
decrease the deformability of the material. This is due to the
compaction of the structure and the redistribution of stresses in the
geocomposite layers. At the same time, the material retains high
chemical resistance and frost resistance, which indicates its durability
and resistance to climatic influences.

An analysis of the correlations between the physico-mechanical
characteristics of the geocomposite showed (Table 2) that most of the
strength parameters, including surface density, tensile and tear
strength in the longitudinal and transverse directions, have a strong
positive correlation with the age of the sample (+0.95...+0.99). This
indicates a tendency for the strength and density of the material to
increase over time. At the same time, elongation at maximum load and at
break in the longitudinal direction shows a strong negative correlation
with age (-0.95...-0.99), which indicates a decrease in plasticity and
increased rigidity of the material. As the service life increases, the
geocomposite becomes more durable but less deformable, which is
associated with aging processes and partial degradation of the polymer
component. Some parameters, such as elongation at the yield point in the
transverse direction and frost resistance, practically do not correlate
with age, which reflects their stability regardless of the operating
time. Positive correlations between strength characteristics and
resistance to aggressive environments (+0.95...+0.96) confirm the
durability of the material under chemically aggressive influences.

Thus, the conducted studies have confirmed that the LPP and
geotextile-based geocomposite has a set of operational properties that
ensure its effective use as a covering material for winter landfill
blocks during long-term operations.

{\bfseries Conclusions.} In the course of the experimental studies, the
physicomechanical properties of a geocomposite based on low-pressure
polyethylene (LPP) and non-woven geotextile used as a heat-protective
coating for winter blocks of TMF dumps of the Kounrad mine were studied.
Based on the analysis of the data obtained, patterns of changes in the
strength and deformation characteristics of the material have been
established depending on the service life. The results obtained made it
possible to evaluate the influence of climatic and operational factors
on the durability and stability of the geocomposite properties:
\end{multicols}

\begin{sidewaystable}
\tcap{Table 2 - Matrix of correlation coefficients between the physical and mechanical characteristics of the geocomposite}
\definecolor{Astral}{rgb}{0.18,0.454,0.709}
\definecolor{RegentStBlue}{rgb}{0.611,0.76,0.898}
\definecolor{TiaMaria}{rgb}{0.768,0.349,0.066}
\definecolor{Spindle}{rgb}{0.741,0.839,0.933}
\definecolor{LinkWater}{rgb}{0.87,0.917,0.964}
\definecolor{PineGlade}{rgb}{0.658,0.815,0.552}
\definecolor{Tacao}{rgb}{0.956,0.69,0.513}
\definecolor{Maize}{rgb}{0.968,0.792,0.674}
\begin{tblr}[
  label = none,
  entry = none,
]{
  width = \linewidth,
  colspec = {Q[5]Q[90]Q[10]Q[10]Q[10]Q[10]Q[10]Q[10]Q[10]Q[10]Q[10]Q[10]Q[10]Q[10]Q[10]Q[10]Q[10]},
  cells = {c},
  cells = {font = \footnotesize},
  row{1} = {font=\bfseries\footnotesize},
  column{1} = {font=\bfseries\footnotesize},
  cell{2}{3} = {Astral},
  cell{2}{4} = {RegentStBlue},
  cell{2}{5} = {RegentStBlue},
  cell{2}{6} = {TiaMaria},
  cell{2}{7} = {TiaMaria},
  cell{2}{8} = {RegentStBlue},
  cell{2}{9} = {RegentStBlue},
  cell{2}{10} = {TiaMaria},
  cell{2}{11} = {Spindle},
  cell{2}{12} = {LinkWater},
  cell{2}{13} = {Spindle},
  cell{2}{14} = {RegentStBlue},
  cell{2}{15} = {PineGlade},
  cell{2}{16} = {RegentStBlue},
  cell{2}{17} = {LinkWater},
  cell{3}{3} = {RegentStBlue},
  cell{3}{4} = {Astral},
  cell{3}{5} = {RegentStBlue},
  cell{3}{6} = {TiaMaria},
  cell{3}{7} = {TiaMaria},
  cell{3}{8} = {RegentStBlue},
  cell{3}{9} = {RegentStBlue},
  cell{3}{10} = {TiaMaria},
  cell{3}{11} = {Spindle},
  cell{3}{12} = {Spindle},
  cell{3}{13} = {RegentStBlue},
  cell{3}{14} = {RegentStBlue},
  cell{3}{15} = {PineGlade},
  cell{3}{16} = {RegentStBlue},
  cell{3}{17} = {LinkWater},
  cell{4}{3} = {RegentStBlue},
  cell{4}{4} = {RegentStBlue},
  cell{4}{5} = {Astral},
  cell{4}{6} = {TiaMaria},
  cell{4}{7} = {TiaMaria},
  cell{4}{8} = {RegentStBlue},
  cell{4}{9} = {RegentStBlue},
  cell{4}{10} = {TiaMaria},
  cell{4}{11} = {LinkWater},
  cell{4}{12} = {Spindle},
  cell{4}{13} = {Spindle},
  cell{4}{14} = {RegentStBlue},
  cell{4}{15} = {PineGlade},
  cell{4}{16} = {RegentStBlue},
  cell{4}{17} = {LinkWater},
  cell{5}{3} = {TiaMaria},
  cell{5}{4} = {TiaMaria},
  cell{5}{5} = {TiaMaria},
  cell{5}{6} = {Astral},
  cell{5}{7} = {RegentStBlue},
  cell{5}{8} = {TiaMaria},
  cell{5}{9} = {TiaMaria},
  cell{5}{10} = {RegentStBlue},
  cell{5}{11} = {Tacao},
  cell{5}{12} = {Maize},
  cell{5}{13} = {Tacao},
  cell{5}{14} = {TiaMaria},
  cell{5}{15} = {PineGlade},
  cell{5}{16} = {TiaMaria},
  cell{5}{17} = {Maize},
  cell{6}{3} = {TiaMaria},
  cell{6}{4} = {TiaMaria},
  cell{6}{5} = {TiaMaria},
  cell{6}{6} = {RegentStBlue},
  cell{6}{7} = {Astral},
  cell{6}{8} = {TiaMaria},
  cell{6}{9} = {TiaMaria},
  cell{6}{10} = {RegentStBlue},
  cell{6}{11} = {Maize},
  cell{6}{12} = {Maize},
  cell{6}{13} = {Tacao},
  cell{6}{14} = {TiaMaria},
  cell{6}{15} = {PineGlade},
  cell{6}{16} = {TiaMaria},
  cell{6}{17} = {Maize},
  cell{7}{3} = {RegentStBlue},
  cell{7}{4} = {RegentStBlue},
  cell{7}{5} = {RegentStBlue},
  cell{7}{6} = {TiaMaria},
  cell{7}{7} = {TiaMaria},
  cell{7}{8} = {Astral},
  cell{7}{9} = {RegentStBlue},
  cell{7}{10} = {TiaMaria},
  cell{7}{11} = {Spindle},
  cell{7}{12} = {LinkWater},
  cell{7}{13} = {Spindle},
  cell{7}{14} = {RegentStBlue},
  cell{7}{15} = {PineGlade},
  cell{7}{16} = {RegentStBlue},
  cell{7}{17} = {LinkWater},
  cell{8}{3} = {RegentStBlue},
  cell{8}{4} = {RegentStBlue},
  cell{8}{5} = {RegentStBlue},
  cell{8}{6} = {TiaMaria},
  cell{8}{7} = {TiaMaria},
  cell{8}{8} = {RegentStBlue},
  cell{8}{9} = {Astral},
  cell{8}{10} = {TiaMaria},
  cell{8}{11} = {LinkWater},
  cell{8}{12} = {LinkWater},
  cell{8}{13} = {Spindle},
  cell{8}{14} = {RegentStBlue},
  cell{8}{15} = {PineGlade},
  cell{8}{16} = {RegentStBlue},
  cell{8}{17} = {LinkWater},
  cell{9}{3} = {TiaMaria},
  cell{9}{4} = {TiaMaria},
  cell{9}{5} = {TiaMaria},
  cell{9}{6} = {RegentStBlue},
  cell{9}{7} = {RegentStBlue},
  cell{9}{8} = {TiaMaria},
  cell{9}{9} = {TiaMaria},
  cell{9}{10} = {Astral},
  cell{9}{11} = {Tacao},
  cell{9}{12} = {Maize},
  cell{9}{13} = {Tacao},
  cell{9}{14} = {TiaMaria},
  cell{9}{15} = {PineGlade},
  cell{9}{16} = {TiaMaria},
  cell{9}{17} = {Maize},
  cell{10}{3} = {Spindle},
  cell{10}{4} = {Spindle},
  cell{10}{5} = {LinkWater},
  cell{10}{6} = {Tacao},
  cell{10}{7} = {Maize},
  cell{10}{8} = {Spindle},
  cell{10}{9} = {LinkWater},
  cell{10}{10} = {Tacao},
  cell{10}{11} = {Astral},
  cell{10}{12} = {LinkWater},
  cell{10}{13} = {LinkWater},
  cell{10}{14} = {Spindle},
  cell{10}{15} = {PineGlade},
  cell{10}{16} = {LinkWater},
  cell{10}{17} = {LinkWater},
  cell{11}{3} = {LinkWater},
  cell{11}{4} = {Spindle},
  cell{11}{5} = {Spindle},
  cell{11}{6} = {Maize},
  cell{11}{7} = {Maize},
  cell{11}{8} = {LinkWater},
  cell{11}{9} = {LinkWater},
  cell{11}{10} = {Maize},
  cell{11}{11} = {LinkWater},
  cell{11}{12} = {Astral},
  cell{11}{13} = {Spindle},
  cell{11}{14} = {Spindle},
  cell{11}{15} = {PineGlade},
  cell{11}{16} = {Spindle},
  cell{11}{17} = {LinkWater},
  cell{12}{3} = {Spindle},
  cell{12}{4} = {RegentStBlue},
  cell{12}{5} = {Spindle},
  cell{12}{6} = {Tacao},
  cell{12}{7} = {Tacao},
  cell{12}{8} = {Spindle},
  cell{12}{9} = {Spindle},
  cell{12}{10} = {Tacao},
  cell{12}{11} = {LinkWater},
  cell{12}{12} = {Spindle},
  cell{12}{13} = {Astral},
  cell{12}{14} = {Spindle},
  cell{12}{15} = {PineGlade},
  cell{12}{16} = {RegentStBlue},
  cell{12}{17} = {LinkWater},
  cell{13}{3} = {RegentStBlue},
  cell{13}{4} = {RegentStBlue},
  cell{13}{5} = {RegentStBlue},
  cell{13}{6} = {TiaMaria},
  cell{13}{7} = {TiaMaria},
  cell{13}{8} = {RegentStBlue},
  cell{13}{9} = {RegentStBlue},
  cell{13}{10} = {TiaMaria},
  cell{13}{11} = {LinkWater},
  cell{13}{12} = {Spindle},
  cell{13}{13} = {Spindle},
  cell{13}{14} = {Astral},
  cell{13}{15} = {PineGlade},
  cell{13}{16} = {RegentStBlue},
  cell{13}{17} = {LinkWater},
  cell{14}{3} = {PineGlade},
  cell{14}{4} = {PineGlade},
  cell{14}{5} = {PineGlade},
  cell{14}{6} = {PineGlade},
  cell{14}{7} = {PineGlade},
  cell{14}{8} = {PineGlade},
  cell{14}{9} = {PineGlade},
  cell{14}{10} = {PineGlade},
  cell{14}{11} = {PineGlade},
  cell{14}{12} = {PineGlade},
  cell{14}{13} = {PineGlade},
  cell{14}{14} = {PineGlade},
  cell{14}{15} = {Astral},
  cell{14}{16} = {PineGlade},
  cell{14}{17} = {PineGlade},
  cell{15}{3} = {RegentStBlue},
  cell{15}{4} = {RegentStBlue},
  cell{15}{5} = {RegentStBlue},
  cell{15}{6} = {TiaMaria},
  cell{15}{7} = {TiaMaria},
  cell{15}{8} = {RegentStBlue},
  cell{15}{9} = {RegentStBlue},
  cell{15}{10} = {TiaMaria},
  cell{15}{11} = {LinkWater},
  cell{15}{12} = {Spindle},
  cell{15}{13} = {RegentStBlue},
  cell{15}{14} = {RegentStBlue},
  cell{15}{15} = {PineGlade},
  cell{15}{16} = {Astral},
  cell{15}{17} = {LinkWater},
  cell{16}{3} = {LinkWater},
  cell{16}{4} = {LinkWater},
  cell{16}{5} = {LinkWater},
  cell{16}{6} = {Maize},
  cell{16}{7} = {Maize},
  cell{16}{8} = {LinkWater},
  cell{16}{9} = {LinkWater},
  cell{16}{10} = {Maize},
  cell{16}{11} = {LinkWater},
  cell{16}{12} = {LinkWater},
  cell{16}{13} = {LinkWater},
  cell{16}{14} = {LinkWater},
  cell{16}{15} = {PineGlade},
  cell{16}{16} = {LinkWater},
  cell{16}{17} = {Astral},
  hlines,
  vlines,
}
№ & Indicator & 1 & 2 & 3 & 4 & 5 & 6 & 7 & 8 & 9 & 10 & 11 & 12 & 13 & 14 & 15\\
1 & Surface density, g/m\textsuperscript{2} & 1 & 0.98 & 0.97 & -0.97 & -0.95 & 0.99 & 0.99 & -0.99 & 0.50 & 0.36 & 0.85 & 0.97 & 0.00 & 0.96 & 0.38\\
2 & Tensile strength (longitudinal direction), kN/m & 0.98 & 1 & 0.99 & -0.97 & -0.93 & 0.96 & 0.97 & -0.96 & 0.50 & 0.60 & 0.90 & 0.95 & 0.00 & 0.95 & 0.40\\
3 & Tensile strength (transverse direction), kN/m & 0.97 & 0.99 & 1 & -0.94 & -0.91 & 0.95 & 0.97 & -0.94 & 0.45 & 0.55 & 0.89 & 0.93 & 0.00 & 0.93 & 0.35\\
4 & Elongation at maximum load (in the longitudinal direction), \% & -0.97 & -0.97 & -0.94 & 1 & 0.92 & -0.95 & -0.93 & 0.99 & -0.50 & -0.36 & -0.85 & -0.97 & 0.00 & -0.95 & -0.38\\
5 & Elongation at maximum load (in the transverse direction), \% & -0.95 & -0.93 & -0.91 & 0.92 & 1 & -0.92 & -0.91 & 0.92 & -0.45 & -0.30 & -0.75 & -0.90 & 0.00 & -0.90 & -0.32\\
6 & Tensile strength (in the longitudinal direction), kN/m & 0.99 & 0.96 & 0.95 & -0.95 & -0.92 & 1 & 0.99 & -0.98 & 0.50 & 0.40 & 0.88 & 0.95 & 0.00 & 0.96 & 0.36\\
7 & Tensile strength (in the transverse direction), kN/m & 0.99 & 0.97 & 0.97 & -0.93 & -0.91 & 0.99 & 1 & -0.96 & 0.48 & 0.38 & 0.87 & 0.94 & 0.00 & 0.95 & 0.35\\
8 & Elongation at break (in the longitudinal direction), \% & -0.99 & -0.96 & -0.94 & 0.99 & 0.92 & -0.98 & -0.96 & 1 & -0.50 & -0.35 & -0.85 & -0.96 & 0.00 & -0.95 & -0.38\\
9 & Elongation at break (in the transverse direction), \% & 0.50 & 0.50 & 0.45 & -0.50 & -0.45 & 0.50 & 0.48 & -0.50 & 1 & 0.10 & 0.20 & 0.50 & 0.00 & 0.45 & 0.12\\
10 & Yield strength (in the longitudinal direction), kN/m & 0.36 & 0.60 & 0.55 & -0.36 & -0.30 & 0.40 & 0.38 & -0.35 & 0.10 & 1 & 0.75 & 0.70 & 0.00 & 0.80 & 0.30\\
11 & Yield strength (in the transverse direction), kN/m & 0.85 & 0.90 & 0.89 & -0.85 & -0.75 & 0.88 & 0.87 & -0.85 & 0.20 & 0.75 & 1 & 0.80 & 0.00 & 0.90 & 0.35\\
12 & Elongation at yield strength (in the longitudinal direction), \% & 0.97 & 0.95 & 0.93 & -0.97 & -0.90 & 0.95 & 0.94 & -0.96 & 0.50 & 0.70 & 0.80 & 1 & 0.00 & 0.95 & 0.38\\
13 & Elongation at yield strength (in the transverse direction), \% & 0.00 & 0.00 & 0.00 & 0.00 & 0.00 & 0.00 & 0.00 & 0.00 & 0.00 & 0.00 & 0.00 & 0.00 & 1 & 0.00 & 0.00\\
14 & Resistance to aggressive environments, \% & 0.96 & 0.95 & 0.93 & -0.95 & -0.90 & 0.96 & 0.95 & -0.95 & 0.45 & 0.80 & 0.90 & 0.95 & 0.00 & 1 & 0.35\\
15 & Resistance to repeated freezing and thawing (frost resistance), \% & 0.38 & 0.40 & 0.35 & -0.38 & -0.32 & 0.36 & 0.35 & -0.38 & 0.12 & 0.30 & 0.35 & 0.38 & 0.00 & 0.35 & 1
\end{tblr}
\end{sidewaystable}

\begin{multicols}{2}
1) During the operation of the geocomposite, an increase in surface
density is observed from 445 to 604 g/m\tsp{2}, which is
associated with the penetration of mineral and organic inclusions into
the structure of the material. An increase in density leads to an
increase in the strength and rigidity of the web but reduces its
ductility and vapor permeability.

2) The tensile strength increases both longitudinally and transversely,
reaching 9.5 and 7.7 kN/m, respectively, which indicates a stabilization
of the structure and increased adhesion between LPP and geotextile
layers during operation.

3) The elongation at maximum load decreases (in the longitudinal
direction from 66 to 44\%, in the transverse direction from 67 to 43\%),
which indicates a decrease in the plasticity of the material due to its
compaction and polymer aging.

4) The tensile strength increases with operation (longitudinally up to
2.49 kN/m, transversely up to 2.95 kN/m), confirming the high strength
of the geocomposite and its ability to maintain the integrity of the
structure under prolonged loads.

5) During operation, anisotropy of deformation properties is manifested:
the longitudinal direction is characterized by greater strength, but
less elasticity, while the transverse direction demonstrates a higher
deformability, which contributes to an even distribution of loads.

6) The yield strength in both directions increases (up to 2.69--3.14
kN/m), which indicates the formation of a stable state of the material
capable of withstanding prolonged static loads without destruction.

7) Geocomposite retains high resistance to aggressive environments ---
up to 96\% after three years of operation, which confirms the chemical
inertia of polyethylene and geotextile fibers, as well as high frost
resistance (92-95\%) during multiple cycles of freezing and thawing.

8) Correlation analysis made it possible to quantify the patterns
identified during mechanical tests.: The increase in the density and
strength of the material is accompanied by a decrease in its
deformability, which indicates the strengthening of the geotextile
structure.

In general, geocomposite has demonstrated complex durability and
stability of properties during long-term operation in a sharply
continental climate. The material retains its protective functions and
mechanical strength, which allows it to be recommended for long-term use
in shelter and protection systems for mining facilities.
\end{multicols}

\begin{center}
{\bfseries Литература}
\end{center}

\begin{refs}
1. Кутепова Н.А., Кутепов Ю.И., Васильева А. Д., Кутепов Ю.Ю. Влияние на
устойчивость отвалов гидроизоляции основания с применением
геосинтетических материалов (геомембран)// ГИАБ. Горный
информационно-аналитический бюллетень. -2025. -№ 1. -C.47-65. DOI
10.25018/0236\_1493\_2025\_1\_0\_47.

2. Wu H., Yao C., Li C., Miao M., Zhong Y., Lu Y., Liu T. Review of
application and innovation of geotextiles in geotechnical engineering //
Materials. -2020. -Vol.13(7):1774. DOI
\href{https://doi.org/10.3390/ma13071774}{10.3390/ma13071774}.

3. Al-Barqawi M., Aqel R., Wayne M., Titi H., Elhajjar R. Polymer
Geogrids: A Review of Material, Design and Structure Relationships //
Materials. -2021. -Vol.14(16):4745. DOI 10.3390/ma14164745.

4. Carneiro J.R., Lopes M.d. L. Weathering of a nonwoven polypropylene
geotextile: Field vs. laboratory exposure // Materials. -2022. -Vol.
15(22):8216. DOI 10.3390/ma15228216.

5. Korolev A., Mishnev M., Zherebtsov D., Vatin N. I., Karelina, M.
Polymers under Load and Heating Deformability: Modelling and Predicting
// Polymers. -2021. -Vol.13(3). -P.428. DOI 10.3390/polym13030428.

6. Scholz P., Falkenhagen J., Wachtendorf V., Brüll R., Simon F.-G.
Investigation on the Durability of a Polypropylene Geotextile under
Artificial Aging Scenarios // Sustainability. -2024. - Vol.16(9): 3559.
DOI 10.3390/su16093559.

7. Markiewicz A., Koda E., Kiraga M., Wrzesiński G., Kozanka K.,
Naliwajko M., Vaverková M.D. Polymeric Products in Erosion Control
Applications: A Review // Polymers. -2024. - Vol.16(17): 2490. DOI
10.3390/polym16172490.

8. Chubarenko B., Domnin D., Simon F-G., Scholz P., Leitsin V., Tovpinets
A., Karmanov K., Esiukova E. Change over Time in the Mechanical
Properties of Geosynthetics Used in Coastal Protection in the
South-Eastern Baltic // Journal of Marine Science and Engineering. -
2023. - Vol.11(1). -P.113. DOI 10.3390/jmse11010113.

9. Vicuña L., Jaramillo-Fierro X., Cuenca P.E., Godoy-Paucar B.,
Inga-Lafebre J.D., Chavez Torres J.L., García J.F., Guaya D., Febres
J.D. Evaluation of the Effectiveness of Geogrids Manufactured from
Recycled Plastics for Slope Stabilization - A Case Study // Polymers.
-2024. - Vol.16(8): 1151. DOI 10.3390/polym16081151.

10. ГОСТ Р 56586-2015. Геомембраны гидроизоляционны е полиэтиленовы е
рулонные. Технические условия/ Стандартинформ. - 2016. - 12 c.

11. ГОСТ Р 50277-92 (ИСО 9864-90). Материалы геотекстильные метод
определения поверхностной плотности / Госстандарт России. - 1992. - 4 c.

12. ГОСТ Р 50300-92. Резцы токарные со сменными режущими пластинам и из
сверхтвердых материалов. Технические условия/ Госстандарт России. -
1994. - 7 c.

13. ГОСТ 11262-2017 (ISO 527-2:2 012). Пластмассы. Метод испытания на
растяжение/ Стандартинформ. - 2018. - 24 c.

14. ГОСТ Р 55035-2012. Дороги автомобильные общего пользования.
Материалы геосинтетические для дорожного строительства. Метод
определения устойчивости к агрессивным средам/ Стандартинформ - 2013. -
12 c.

15. ГОСТ Р 55032-2012. Национальный стандарт российской федерации.
Дороги автомобильные общего пользования. Материалы геосинтетические для
дорожного строительства. Метод определения устойчивости к многократному
замораживанию и оттаиванию/ Стандартинформ .- 2013. - 12 c.

16. Carneiro J.R., Gomes M., Vieira C.S. Inclined Plane Shear Behaviour
of a Recycled Construction and Demolition Material--Geocomposite
Interface // International Journal of Geosynthetics and Ground
Engineering. - 2023. - Vol.9: 88. DOI 10.1007/s40891-023-00507-1.

17. Carneiro J.R., Almeida F., Carvalho F., Lopes M.d.L. Tensile and
Tearing Properties of a Geocomposite Mechanically Damaged by Repeated
Loading and Abrasion // Materials. -2023. - Vol.16(21):7047. DOI
10.3390/ma16217047.

18. Ferreira F.B., Vieira C.S., Lopes M.L. Pullout Behavior of Different
Geosynthetics - Influence of Soil Density and Moisture Content //
Frontiers in Built Environment. -2020. - Vol.6: 12. -P.1-13. DOI
10.3389/fbuil.2020.00012.

19. Lesov K., Uralov A., Kenjaliyev M., Ergashev U., Huaiping F. The
importance of using geosynthetic materials in ensuring anti-erosion
stability of railway embankments // Vibroengineering Procedia. -2025.
-Vol.60. -P.318--324. DOI 10.21595/vp.2025.25326.

20. Kim H., Choi J. Influence of crystalline structure on creep
resistance capability in semi-crystalline Polymers: A coarse-grained
molecular dynamics study // International Journal of Fatigue. -2024. -
Vol.188: 108517. DOI 10.1016/j.ijfatigue.2024.108517.

21. Yan W., Gao S., Feng Y., Liu B. Toward excellent strength, modulus,
and high-temperature creep resistance for wide molecular weight
distribution polyethylene fiber by melt-spinning // European Polymer
Journal. -2025. -Vol.236: 114132. DOI 10.1016/j.eurpolymj.2025.114132.

22. Lomovskoy V. A., Shatokhina S.A. Relaxation Phenomena in Low-Density
and High-Density Polyethylene // Polymers. -2024. - Vol.16(24): 3510.
DOI 10.3390/polym16243510.
\end{refs}

\begin{center}
{\bfseries References}
\end{center}

\begin{refs}
1. Kutepova N.A., Kutepov Ju.I., Vasil' eva A. D.,
Kutepov Ju.Ju. Vlijanie na ustojchivost'{} otvalov
gidroizoljacii osnovanija s primeneniem geosinteticheskih materialov
(geomembran)// GIAB. Gornyj informacionno-analiticheskij
bjulleten'. -2025. -№ 1. -C.47-65.
DOI 10.25018/0236\_1493\_2025\_1\_0\_47. {[}in Russian{]}

2. Wu H., Yao C., Li C., Miao M., Zhong Y., Lu Y., Liu T. Review of
application and innovation of geotextiles in geotechnical engineering //
Materials. -2020. -Vol.13(7):1774. DOI
\href{https://doi.org/10.3390/ma13071774}{10.3390/ma13071774}.

3. Al-Barqawi M., Aqel R., Wayne M., Titi H., Elhajjar R. Polymer
Geogrids: A Review of Material, Design and Structure Relationships //
Materials. -2021. -Vol.14(16):4745. DOI 10.3390/ma14164745.

4. Carneiro J.R., Lopes M.d. L. Weathering of a nonwoven polypropylene
geotextile: Field vs. laboratory exposure // Materials. -2022. -Vol.
15(22):8216. DOI 10.3390/ma15228216.

5. Korolev A., Mishnev M., Zherebtsov D., Vatin N. I., Karelina, M.
Polymers under Load and Heating Deformability: Modelling and Predicting
// Polymers. -2021. -Vol.13(3). -P.428. DOI 10.3390/polym13030428.

6. cholz P., Falkenhagen J., Wachtendorf V., Brüll R., Simon F.-G.
Investigation on the Durability of a Polypropylene Geotextile under
Artificial Aging Scenarios // Sustainability. -2024. - Vol.16(9): 3559.
DOI 10.3390/su16093559.

7. Markiewicz A., Koda E., Kiraga M., Wrzesiński G., Kozanka K.,
Naliwajko M., Vaverková M.D. Polymeric Products in Erosion Control
Applications: A Review // Polymers. -2024. - Vol.16(17): 2490. DOI
10.3390/polym16172490.

8. Chubarenko B., Domnin D., Simon F-G., Scholz P., Leitsin V.,
Tovpinets A., Karmanov K., Esiukova E. Change over Time in the
Mechanical Properties of Geosynthetics Used in Coastal Protection in the
South-Eastern Baltic // Journal of Marine Science and Engineering. -
2023. - Vol.11(1). -P.113. DOI 10.3390/jmse11010113.

9. Vicuña L., Jaramillo-Fierro X., Cuenca P.E., Godoy-Paucar B.,
Inga-Lafebre J.D., Chavez Torres J.L., García J.F., Guaya D., Febres
J.D. Evaluation of the Effectiveness of Geogrids Manufactured from
Recycled Plastics for Slope Stabilization - A Case Study // Polymers.
-2024. - Vol.16(8): 1151. DOI 10.3390/polym16081151.

10. GOST R 56586-2015. Geomembrany gidroizoljacionny e polijetilenovy e
rulonnye. Tehnicheskie uslovija/ Standartinform. - 2016. - 12 c. {[}in
Russian{]}

11. GOST R 50277-92 (ISO 9864-90). Materialy
geotekstil' nye metod opredelenija poverhnostnoj
plotnosti / Gosstandart Rossii. - 1992. - 4 c. {[}in Russian{]}

12. GOST R 50300-92. Rezcy tokarnye so smennymi rezhushhimi plastinam i
iz sverhtverdyh materialov. Tehnicheskie uslovija/ Gosstandart Rossii. -
1994. - 7 c. {[}in Russian{]}

13. GOST 11262-2017 (ISO 527-2:2012). Plastmassy. Metod ispytanija na
rastjazhenie/ Standartinform. - 2018. - 24 c. {[}in Russian{]}

14. GOST R 55035-2012. Dorogi avtomobil' nye obshhego
pol' zovanija. Materialy geosinteticheskie dlja
dorozhnogo stroitel' stva. Metod opredelenija
ustojchivosti k agressivnym sredam/ Standartinform - 2013. - 12 c. {[}in
Russian{]}

15. GOST R 55032-2012. Nacional' nyj standart rossijskoj
federacii. Dorogi avtomobil' nye obshhego
pol' zovanija. Materialy geosinteticheskie dlja
dorozhnogo stroitel' stva. Metod opredelenija
ustojchivosti k mnogokratnomu zamorazhivaniju i ottaivaniju/
Standartinform .- 2013. - 12 c. {[}in Russian{]}

16. Carneiro J.R., Gomes M., Vieira C.S. Inclined Plane Shear Behaviour
of a Recycled Construction and Demolition Material--Geocomposite
Interface // International Journal of Geosynthetics and Ground
Engineering. - 2023. - Vol.9: 88. DOI 10.1007/s40891-023-00507-1.

17. Carneiro J.R., Almeida F., Carvalho F., Lopes M.d.L. Tensile and
Tearing Properties of a Geocomposite Mechanically Damaged by Repeated
Loading and Abrasion // Materials. -2023. - Vol.16(21):7047. DOI
10.3390/ma16217047.

18. Ferreira F.B., Vieira C.S., Lopes M.L. Pullout Behavior of Different
Geosynthetics - Influence of Soil Density and Moisture Content //
Frontiers in Built Environment. -2020. - Vol.6: 12. -P.1-13. DOI
10.3389/fbuil.2020.00012.

19. Lesov K., Uralov A., Kenjaliyev M., Ergashev U., Huaiping F. The
importance of using geosynthetic materials in ensuring anti-erosion
stability of railway embankments // Vibroengineering Procedia. -2025.
-Vol.60. -P.318--324. DOI 10.21595/vp.2025.25326.

20. Kim H., Choi J. Influence of crystalline structure on creep
resistance capability in semi-crystalline Polymers: A coarse-grained
molecular dynamics study // International Journal of Fatigue. -2024. -
Vol.188: 108517. DOI 10.1016/j.ijfatigue.2024.108517.

21. Yan W., Gao S., Feng Y., Liu B. Toward excellent strength, modulus,
and high-temperature creep resistance for wide molecular weight
distribution polyethylene fiber by melt-spinning // European Polymer
Journal. -2025. -Vol.236: 114132. DOI 10.1016/j.eurpolymj.2025.114132.

22. Lomovskoy V. A., Shatokhina S.A. Relaxation Phenomena in Low-Density
and High-Density Polyethylene // Polymers. -2024. - Vol.16(24): 3510.
DOI 10.3390/polym16243510.
\end{refs}

\begin{info}
\hspace{1em}\emph{{\bfseries Information about the authors}}

Seraya N.V. - Candidate of Chemical Sciences, Professor, D. Serikbayev
East Kazakhstan Technical University, Ust-Kamenogorsk, Kazakhstan,
e-mail: nseraya@mail.ru;

Litvinov V.V. - Deputy Director, Proektno-ekologicheskoe bjuro LLP,
Ust-Kamenogorsk, Kazakhstan, e-mail: litvinov\_vadim@mail.ru;

Daumova G.K. - Candidate of Technical Sciences, Professor,
D.Serikbayev East Kazakhstan Technical University, Ust-Kamenogorsk,
Kazakhstan, е-mail: gulzhan.daumova@mail.ru;

Sovetkhanov D.E. - Master of Technical Sciences, Doctoral student, D.
Serikbayev East Kazakhstan Technical University, Ust-Kamenogorsk,
Kazakhstan, е-mail: sovetkhanovdauren5@gmail.com;

Rutkowska-Gorczyca M. - ~~PhD, Professor, Wrocław University of
Science and Technology, Wrocław, Poland, е-mail:
malgorzata.rutkowska-gorczyca@pwr.edu.pl;

Baimurzina M.M. - student, D. Serikbayev East Kazakhstan Technical
University, Ust-Kamenogorsk, Kazakhstan; е-mail: may787585@gmail.com.

\hspace{1em}\emph{{\bfseries Сведения об авторах}}

Серая Н.В. - кандидат химических наук, профессор,
Восточно-Казахстанский технический университет им. Д.Серикбаева,
Усть-Каменогорск, Казахстан, е-mail: nseraya@mail.ru;

Литвинов В.В. - заместитель директора, ТОО «Проектно-экологическое
бюро», Усть-Каменогорск, Казахстан, е-mail: litvinov\_vadim@mail.ru;

Даумова Г. К. -- кандидат технических наук, профессор,
Восточно-Казахстанский технический университет им. Д.Серикбаева,
Усть-Каменогорск, Казахстан, е-mail: gulzhan.daumova@mail.ru;

Советханов Д.Е. -- магистр технических наук, докторант,
Восточно-Казахстанский технический университет им. Д.Серикбаева,
Усть-Каменогорск, Казахстан, e-mail: sovetkhanovdauren5@gmail.com;

Rutkowska-Gorczyca M.- PhD, профессор, Вроцлавский университет науки и
технологии, Вроцлав, Польша, е-mail:
malgorzata.rutkowska-gorczyca@pwr.edu.pl.;

Баймурзина М.М. - студент, Восточно-Казахстанский технический
университет им. Д.Серикбаева, Усть-Каменогорск, Казахстан, e-mail:
may787585@gmail.com.
\end{info}
